% answers_short.tex —— 2008 年·普通高等学校招生考试·重庆卷·理科 · 简短答案速查
% 由 tools/make_short_answers.py 自 answers.tex 的 \gksolhead/\gkanswer 提取，
% 并与 2008/2008_数学真题_重庆理_答案要点.md 逐题比对后生成，勿手改（改源 MD 后重跑）。

\begin{gknotice}
\textbf{使用说明}：本节只给最终答案便于快速核对，完整推导见解析版。答案为编写者独立推导，\textbf{非官方标准答案}。
\end{gknotice}

\gkpart{选择题}
\begin{center}\begin{tabular}{ccccccccccc}
\toprule
题号 & 1 & 2 & 3 & 4 & 5 & 6 & 7 & 8 & 9 & 10 \\
\midrule
答案 & A & A & B & C & D & C & A & C & D & B \\
\bottomrule
\end{tabular}\end{center}

\gkpart{填空题}
\begin{itemize}[leftmargin=2.6em]
  \item[\textbf{11.}] $\{2,5\}$。
  \item[\textbf{12.}] $\dfrac13$。
  \item[\textbf{13.}] $3$。
  \item[\textbf{14.}] $-72$。
  \item[\textbf{15.}] $y=x+1$。
  \item[\textbf{16.}] $216$。
\end{itemize}

\gkpart{解答题}
\begin{itemize}[leftmargin=2.6em]
  \item[\textbf{17.}] （1）$\dfrac{\sqrt7}3$；（2）$\dfrac{14\sqrt3}9$。
  \item[\textbf{18.}] （1）$\dfrac14$；（2）分布列见解析，$E\xi=\dfrac{47}{16}$。
  \item[\textbf{19.}] （1）$2$；（2）$\arccos\dfrac{3\sqrt{34}}{34}$。
  \item[\textbf{20.}] （1）$b=2a$，$c=2a+3$；（2）增区间 $(-2,2)$，减区间 $(-\infty,-2)$、$(2,+\infty)$。
  \item[\textbf{21.}] （1）$\dfrac{x^2}9+\dfrac{y^2}5=1$；（2）$P\left(\dfrac{3\sqrt3}2,\dfrac{\sqrt5}2\right)$、$\left(\dfrac{3\sqrt3}2,-\dfrac{\sqrt5}2\right)$、$\left(-\dfrac{3\sqrt3}2,\dfrac{\sqrt5}2\right)$、$\left(-\dfrac{3\sqrt3}2,-\dfrac{\sqrt5}2\right)$。
  \item[\textbf{22.}] （1）$a_3=16$，$a_4=\dfrac1{256}$，$a_{2008}=2^{-2^{2007}}$；（2）$a_2=\sqrt2$，$b_n=2^{2-2^{1-n}}$。
\end{itemize}
